% This is the main file for the ulm class diagram for the robotic path planning project
\documentclass[10pt]{article}
% Standard packages
\usepackage[legalpaper, landscape, top= 0.0in, bottom=0.5in, right=0.5in, left=0.5in]{geometry}
\usepackage[active,tightpage]{preview}
\renewcommand{\PreviewBorder}{0.5in}
% Required packages for tikz-uml
\usepackage{tikz}
\usepackage{ifthen}
\usepackage{xstring}
\usepackage{calc}
\usepackage{pgfplots}
\usepackage{tikz-uml}
\newcommand{\dunder}{\rule{0.1in}{0.1pt}}

\title{UML Class Diagram for Robotic Path Planning}
\author{Bruun, E. \and Flores, J. \and Kumar, V.}

\declare{}
\begin{document}
\maketitle

\begin{tikzpicture}
\begin{umlpackage}{RRT}

%% Input and Parameter Classes %%

%Parameters (ABSTRACT)
\setvalue{Xcoord=16.}
\setvalue{Ycoord=2.}
\setvalue{aligninfo=north}
\setvalue{widthinfo=6cm}
\newcommand{\relstr}{}
%This class provides a standard interface to an input file which provides at run-time for program parameters such as step_size
\umlclass[x=\getvalue{Xcoord}, y=\getvalue{Ycoord},
				type = abstract,
				\relstr,
				anchor=\getvalue{aligninfo},
				width = \getvalue{widthinfo}]
				{Params}
% Data
{
+ file$\_$name: string
}
% Methods
{
\dunder{init}\dunder(filename)\\
+ parse$\_$parameters(string file$\_$name) \\
+ void declare$\_$entry(string entry$\_$name, \\ string default$\_$value, \\ Reg$\_$Ex(?) pattern)\\
+ set(string entry$\_$name, string new$\_$value)\\
+ print$\_$paramters()
}
	

%Domain Parameters
\setvalue{aligninfo=north}
\setvalue{widthinfo=5cm}
\renewcommand{\relstr}{below right= 0cm and 9cm of Params.north}
%This class provides a standard interface to an input file which provides at run-time for program parameters such as step_size
\umlclass[x=\getvalue{Xcoord}, y=\getvalue{Ycoord},
				\relstr,
				anchor=\getvalue{aligninfo},
				width = \getvalue{widthinfo}]
				{ParamsDomain}
% Data
{
+ file$\_$name: string \\
+ info$\_$dims: string \\
+ info$\_$domain: string \\
+ info$\_$obstacles: string
}
% Methods
{
}
	
\umlinherit[anchors= 180 and 0, geometry=-|-]{ParamsDomain}{Params}

%Origin/Goal Parameters
\setvalue{aligninfo=north}
\setvalue{widthinfo=5cm}
\renewcommand{\relstr}{below right=3cm and 9cm of Params.north}
%This class provides a standard interface to an input file which provides at run-time for program parameters such as step_size
\umlclass[x=\getvalue{Xcoord}, y=\getvalue{Ycoord},
				\relstr,
				anchor = \getvalue{aligninfo},
				width = \getvalue{widthinfo}]
				{ParamsOriginGoal}
% Data
{
+ file$\_$name: string \\
+ info$\_$origin: string \\
+ info$\_$goal: string \\
}
% Methods
{
}
	
\umlinherit[anchors= 180 and 0, geometry=-|-]{ParamsOriginGoal}{Params}

%RRT Parameters
\setvalue{aligninfo=north}
\setvalue{widthinfo=5cm}
\renewcommand{\relstr}{below right=6cm and 9cm of Params.north}
%This class provides a standard interface to an input file which provides at run-time for program parameters such as step_size
\umlclass[x=\getvalue{Xcoord}, y=\getvalue{Ycoord},
				\relstr,
				anchor=\getvalue{aligninfo},
				width = \getvalue{widthinfo}]
				{ParamsRRT}
% Data
{
+ ....tbd
}
% Methods
{
}
	
\umlinherit[anchors= 180 and 0, geometry=-|-]{ParamsRRT}{Params}



%% Geometry Classes %%

%Shape (ABSTRACT)
\setvalue{Xcoord=7.}
\setvalue{Ycoord=2}
\setvalue{aligninfo=north}
\setvalue{widthinfo=4cm}
\renewcommand{\relstr}{}
\umlclass[x=\getvalue{Xcoord}, y=\getvalue{Ycoord}, type=abstract]{Shape}
% Data
{
+ dim: int \\
+ shape$\_$name: string
}
% Methods
{
+  \textit{inside(Point p)} \\
+  \textit{intersects(Point p1, Point p2)}
}

%Rectangle
\setvalue{aligninfo=north}
\setvalue{widthinfo=4cm}
\renewcommand{\relstr}{below left=0cm and 7 cm of Shape.north}
\umlclass[x=\getvalue{Xcoord}, y=\getvalue{Ycoord}]{Rectangle}
% Data
{
+ dim: int \\
+ bottom$\_$left: Point \\
+ top$\_$right: Point  \\
+ shape$\_$name: string
}
% Methods
{
+ \dunder{init}\dunder (Point p1, Point p2) \\
+ inside(Point p)\\
+ intersects(Point p1, Point p2)
}

\umlinherit[anchors= 0 and 180, geometry=-|-]{Rectangle}{Shape}

%Circle
\setvalue{aligninfo=north}
\setvalue{widthinfo=4cm}
\renewcommand{\relstr}{below left=4cm and 7 cm of Shape.north}
\umlclass[x=0, y=-7]{Circle}
% Data
{
+ dim: int \\
+ radius: double \\
+ center: Point \\
+ shape$\_$name: string
}
% Methods
{
+ \dunder{init}\dunder(double radius, Point center) \\
+ inside(Point p)\\
+ intersects(Point p1, Point p2)
}

\umlinherit[anchors= 0 and 180, geometry=-|-]{Circle}{Shape}

%Freeform
\setvalue{aligninfo=north}
\setvalue{widthinfo=4cm}
\renewcommand{\relstr}{below left=8cm and 7 cm of Shape.north}
\umlclass[x=\getvalue{Xcoord}, y=\getvalue{Ycoord}]{Freeform}
% Data
{
+ dim: 2 \\
+ shape$\_$name: string \\
+ perim$\_$step : double \\
+ perim$\_$pnts : ndarray
+ image
}
% Methods
{
+ \dunder{init}\dunder(filename) \\
+ inside(pt) : bool \\
+ intersects(pt1, pt2) : bool \\
+ get$\_$image(filename) \\
+ get$\_$perim() : ndarray
}

\umlinherit[anchors= 0 and 180, geometry=-|-]{Freeform}{Shape}


%% Domain Classes %%

%Domain
\setvalue{Xcoord=13.}
\setvalue{Ycoord=-5.}
\setvalue{aligninfo=north}
\setvalue{widthinfo=4cm}
\renewcommand{\relstr}{}
\umlclass[x=10, y=-2]{Domain}
% Data
{
+ dim: int \\
+ points: \\
+ Obstacles: \\
+ Origin: \\
+ Goal: \\
}
% Methods
{
CreateDomain(): DO
}

%Obstacles
\setvalue{aligninfo=north}
\setvalue{widthinfo=5cm}
\renewcommand{\relstr}{below left=-2cm and 6 cm of Domain.north}
\umlclass[x=\getvalue{Xcoord}, y=\getvalue{Ycoord}]{Obstacle}
% Data
{
+ dim: int \\
+ shape$\_$name: string\\
+ 
}
% Methods
{
+ \dunder{init}\dunder(shape$\_$name)
}
\umlaggreg[anchors= 180 and 0, geometry=-|-]{Domain}{Obstacle}

%Origin
\setvalue{aligninfo=north}
\setvalue{widthinfo=5cm}
\renewcommand{\relstr}{below left=1cm and 6 cm of Domain.north}
\umlclass[x=\getvalue{Xcoord}, y=\getvalue{Ycoord},
				\relstr,
				anchor=\getvalue{aligninfo},
				width = \getvalue{widthinfo}]
				{Origin}
% Data
{
+ shape$\_$name: string\\
}
% Methods
{
+ \dunder{init}\dunder(shape$\_$name)
}
\umlaggreg[anchors= 180 and 0, geometry=-|-]{Domain}{Origin}

%Point
\setvalue{aligninfo=north}
\setvalue{widthinfo=5cm}
\renewcommand{\relstr}{below left=4cm and 6 cm of Domain.north}
\umlclass[x=\getvalue{Xcoord}, y=\getvalue{Ycoord},
				\relstr,
				anchor=\getvalue{aligninfo},
				width = \getvalue{widthinfo}]
				{Goal}
% Data
{
+ shape$\_$name: string\\
}
% Methods
{
+ \dunder{init}\dunder(shape$\_$name)
}
\umlaggreg[anchors= 180 and 0, geometry=-|-]{Domain}{Goal}



%% Algorithm Classes %%

%Solver
\setvalue{Xcoord=10.}
\setvalue{Ycoord=-13.}
\setvalue{aligninfo=north}
\setvalue{widthinfo=4cm}
\renewcommand{\relstr}{}
\umlclass[x=0, y=-7]{Solver}
% Data
{
+ vertexList: ndarray \\
+ parentList : int list \\
+ costList : float list
}
% Methods
{
+ Solver() \\
+ Solver(dim : int, trials : int) \\
------\\
+ FindRRTPath (method : str, int : trials) : void
}
\umlaggreg[anchors= 90 and -90, geometry=|-|]{Solver}{Domain}

%RRT (ABSTRACT)
\setvalue{Xcoord=10.}
\setvalue{Ycoord=-21.}
\setvalue{aligninfo=north}
\setvalue{widthinfo=4cm}
\renewcommand{\relstr}{}
\umlclass[type = abstract, x=0, y=-15]{RRT}
% Data
{

}
% Methods
{
 \textit{+ step() : void} \\
 
 ------\\
 
 + NewPoint() : Point \\
 + NewBranch() : Branch \\

------\\

 + UpdateList() : void

}
\umlcompo[mult1 = 1, mult2 = 1]{Solver}{RRT}

%Vertex
\setvalue{aligninfo=north}
\setvalue{widthinfo=6.5cm}
\renewcommand{\relstr}{below left=-2.5cm and 8 cm of RRT.north}
\umlclass[x=\getvalue{Xcoord}, y=\getvalue{Ycoord}]{Vertex}
% Data
{
	coord : list float
}
% Methods
{
+ \dunder{init}\dunder (self, dim : int, domain : ndarray) \\ % Initialize Point class by making random coordinates

------\\
+ p$\_$nearest(self, VertexList : ndarray) : int\\ % Nearest point in existing vertex list
+ p$\_$extend(self, stepSize : float) : coord \\% Create new vertex based on max step size
}

\umlaggreg[anchors= 180 and 0, geometry=-|-]{RRT}{Vertex}

%Edge
\setvalue{aligninfo=north}
\setvalue{widthinfo=6.5cm}
\renewcommand{\relstr}{below left=2.5cm and 8 cm of RRT.north}
\umlclass[x=\getvalue{Xcoord}, y=\getvalue{Ycoord}]{Edge}
% Data
{
	start : Point \\
	end : Point
}
% Methods
{
+ \dunder{init}\dunder (start, end) \\

------\\

+ edge$\_$cost(self) : float\\
+ tree$\_$cost(self) : float
}
\umlaggreg[anchors= 180 and 0, geometry=-|-]{RRT}{Edge}

%RRT Regular
\setvalue{aligninfo=north}
\setvalue{widthinfo=6cm}
\renewcommand{\relstr}{below left=5cm and 0 cm of RRT.north}
\umlclass[x=\getvalue{Xcoord}, y=\getvalue{Ycoord}]{RRTRegular}
% Data
{

}
% Methods
{
 + step() : void \\
 
 ------\\
 
 + CheckGoal(Point, goal : ndarray) : bool
}
\umlinherit[anchors= 90 and -90, geometry=|-|]{RRTRegular}{RRT}

%RRT Regular Heuristic
\setvalue{aligninfo=north}
\setvalue{widthinfo=6cm}
\renewcommand{\relstr}{below left=11cm and 0 cm of RRT.north}
\umlclass[x=-1, y=-21]{RRTRegularHeuristic}
% Data
{

}
% Methods
{
Implement if time...
 
}
\umlinherit[anchors= 90 and -90, geometry=|-|]{RRTRegularHeuristic}{RRTRegular}

%RRT Star
\setvalue{aligninfo=north}
\setvalue{widthinfo=6cm}
\renewcommand{\relstr}{below left=5cm and -8 cm of RRT.north}
\umlclass[x=8, y=-18]{RRTStar}
% Data
{
+ searchRadius : float\\
+ neighbourCostList : [int, float, float]\\
}
% Methods
{
 + step() : void \\
 
 ------\\
 
 + FindNearestNeightbours(self) : void\\
 + NeighbourCosts(self) : void \\ %Get total branch costs for neighbours, and also the branch costs to new point
 + CheapestTotalCost(self) : int \\
 + RewireTree() : void \\

}
\umlinherit[anchors= 90 and -90, geometry=|-|]{RRTStar}{RRT}

%RRT Star FN
\setvalue{aligninfo=north}
\setvalue{widthinfo=6cm}
\renewcommand{\relstr}{below left=11cm and -8 cm of RRT.north}
\umlclass[x=8, y=-23]{RRTStarFN}
% Data
{
+ maxNodes : int
}
% Methods
{
-FindSmallBranch(self) : void \\
-DeleteBranch(self) : void 
}
\umlinherit[anchors= 90 and -90, geometry=|-|]{RRTStarFN}{RRTStar}


%% Output Classes %%

%Plot Output
\setvalue{aligninfo=north}
\setvalue{widthinfo=7cm}
\renewcommand{\relstr}{below left=-8cm and -12 cm of Solver.north}
\umlclass[x=\getvalue{Xcoord}, y=\getvalue{Ycoord},
				\relstr,
				anchor=\getvalue{aligninfo},
				width = \getvalue{widthinfo}]
				{OutputPlot}
% Data
{
+ filename: string \\
+ animate: bool \\
+ dimension: int \\
+ data: list \\ % lists scatter and objects to plot
+ layout: (title: string, showlegend: bool) \\  
+ fig: (data, layout) 
}
% Methods
{
+ \dunder{init}\dunder (self, filename: string) \\ % Initialize Plot class with filename 
------\\
+ plot\_vertex (Solver) : void \\
+ plot\_path (Solver) : void \\
+ plot\_obstacles (Input) : void \\
+ plot\_start (Input) : void \\
+ plot\_goal (Input) : void \\
+ draw (animate) : void % Plot to the file specified  
}
\umlaggreg[anchors= 180 and 0, geometry=-|-]{OutputPlot}{Solver}
\umlaggreg[anchors= 90 and 0, geometry=|-]{OutputPlot}{Domain}

%Text Output
\setvalue{aligninfo=north}
\setvalue{widthinfo=7cm}
\renewcommand{\relstr}{below left=-0cm and -12 cm of Solver.north}
\umlclass[x=\getvalue{Xcoord}, y=\getvalue{Ycoord},
				\relstr,
				anchor=\getvalue{aligninfo},
				width = \getvalue{widthinfo}]
				{OutputText}
% Data
{
+ output$\_$file: string \\
}
% Methods
{
+ \dunder{init}\dunder () \\
+ format$\_$data() \\
+ write$\_$output()
}
\umlaggreg[anchors= 180 and 0, geometry=-|-]{OutputText}{Solver}

\end{umlpackage}

\end{tikzpicture}

\end{document}
